\documentclass[multi=tikzpicture,border=9pt]{standalone}
\usepackage[UTF8]{ctex}
\usepackage{amsmath,amssymb}
\usepackage{tikz}
\usetikzlibrary{arrows.meta,calc,patterns}
\definecolor{main}{HTML}{17365D}
\definecolor{accent}{HTML}{AC3931}
\tikzset{every picture/.style={font=\small,line width=.75pt},>={Stealth[length=2mm]},arr/.style={->},box/.style={draw=main,rounded corners=2pt,fill=main!4,text width=3.6cm,minimum height=1cm,align=center},note/.style={align=center,text width=9cm}}
\begin{document}

\begin{tikzpicture}
\draw[main,thick] (0,0) circle (2);\fill[accent] (0,0) circle (3pt);\node[below] at (0,0){核 $+q$};\fill[main] (2,0) circle (3pt);\node[right] at (2,0){电子 $-q$};\draw[arr,accent] (2,0)--(2,1.2) node[right]{$\mathbf v$};\draw[dashed] (0,0)--(2,0) node[midway,below]{$r_0$};\node at (-.7,.7){$B_0\otimes$};\node[note] at (0,-3.3){电子逆时针，等效电流顺时针，核处自生磁场向纸内};
\end{tikzpicture}

\begin{tikzpicture}
\draw[main,thick] (0,0) circle (2);\draw[dashed] (0,0)--(2,0);\fill (2,0) circle (3pt);\draw[arr,main] (2,0)--(3.2,0) node[right]{$\dot r$};\draw[arr,accent] (2,0)--(2,1.4) node[right]{$F_\varphi=qB\dot r$};\node at (0,.35){$B\odot$};\node[note] at (0,-3.2){外场向纸外，径向速度产生切向磁力；$\dot L_z=qBr\dot r$};
\end{tikzpicture}

\begin{tikzpicture}
\node[box] (n0) at (0,0) {有向机械量\\$L_z=mr^2\dot\varphi$};
\node[box] (n1) at (5,0) {轴向守恒量\\$J=L_z-qBr^2/2$};
\node[box] (n2) at (5,-2.3) {展开动能\\保留 $qBJ/(2m)$};
\node[box] (n3) at (0,-2.3) {能级到一阶\\$E_{n,s}=E_n^0+sn\mu_BB$};
\draw[arr] (n0)--(n1);\draw[arr] (n1)--(n2);\draw[arr] (n2)--(n3);
\node[note] at (2.5,-4.1) {求半径时常量项导数为零，求能量时不能删除};
\end{tikzpicture}

\begin{tikzpicture}
\draw[main,thick] (0,0)--(2,0) node[right]{$E_1^{(0)}$};\draw[main,thick] (0,3)--(2,3) node[right]{$E_2^{(0)}$};\draw[accent,thick] (4,-.25)--(6,-.25);\draw[accent,thick] (4,.25)--(6,.25);\draw[accent,thick] (4,2.5)--(6,2.5);\draw[accent,thick] (4,3.5)--(6,3.5);\draw[arr] (4.5,3.5)--(4.5,.25);\draw[arr] (5.5,2.5)--(5.5,-.25);\node[below] at (1,-.8){无场};\node[below] at (5,-.8){弱场示意};\node[note] at (3,-2){教学圆轨道模型的同绕行方向分支\\$\Delta\nu=\pm\mu_B B/h$；分裂尺度夸大以示意};
\end{tikzpicture}

\begin{tikzpicture}
\node[box] at (0,0){半径方程\\$J^2=mk r+q^2B^2r^4/4$\\首修正为 $B^2$};\node[box] at (5,0){能量展开\\$qBJ/(2m)$\\已有 $B$ 的首修正};\node[note] at (2.5,-2.5){半径无一阶变化，不等于光谱无一阶变化；保持近似精度一致};
\end{tikzpicture}
\end{document}
